% !TeX program = lualatex
% !TeX encoding = utf8
% !TeX spellcheck = uk_UA
% !TeX root =../ClassicalEletrodynamics.tex

%=========================================================
\chapter*{Передмова}\label{\currfilebase}
\addcontentsline{toc}{chapter}{Передмова}
%=========================================================


Посібник містить матеріали з класичної електродинаміки в рамках програми  з теоретичної фізики, розрахованої на студентів третього курсу навчально-наукового фізико-технічного факультету КПІ імені Ігоря Сікорського.


Викладено основні поняття класичної електродинаміки, розглянуто рівняння Максвелла, їх властивості, основні розв'язки, процеси випромінювання та розсіювання, вільні поля та їх характеристики. Посібник складається з двох частин:
І. Класична електродинаміка; ІІ. Чотиривимірне формулювання класичної електродинаміки.

У частині І викладено базові поняття та рівняння, розв'язки рівнянь Максвелла, що описують ізольовану систему зарядів і струмів потенціали ізольованої системи зарядів і струмів, задача Коші для рівнянь Максвелла). Обговорено розв'язки для вільного електромагнітного поля та процеси випромінювання.

У частині ІІ викладено базові принципи та поняття спеціальної теорії відносності (СТВ) близько до оригінальної роботи Айнштайна. Співвідношення СТВ переписано у просторі Мінковського.  Далі рівняння Максвелла та основні пов'язані з ними величини переписано у коваріантній чотиривимірній формі за допомогою тензора електромагнітного поля та чотиривектора густини струму та розглянуто основні наслідки. Виведено коваріантні рівняння руху заряджених частинок в електромагнітному полі за допомогою принципу відповідності з класичною динамікою. Розглянуто варіаційний принцип для рівнянь електродинаміки та закони збереження.

Посібник містить вправи для кращого засвоєння матеріалу.
%Книга містить чотири додатки: Д.1. Узагальнені функції, Д.2. Застосування спеціальних функцій для розв’язання деяких задач  електродинаміки, Д.3. Контрольні запитання. Д.4. Приклади задач для самостійного розв’язання.
